\documentclass[11pt]{article}
\usepackage{amsmath,amssymb,amsthm,hyperref}
\newtheorem{theorem}{Theorem}
\newtheorem{lemma}[theorem]{Lemma}
\begin{document}
\title{Deep Real Memory and the Cost of a Flat Rational Decoder}
\author{}
\date{November 24, 2026}
\maketitle

\paragraph{Status and scope.}
This is a primary-source transfer audit with complete written proofs
of the stated scalar-decoder obstructions. It contains no Lean proof
and gives no new unrestricted head-complexity lower bound. In
particular, an arbitrary reparameterization of all head weights,
or a joint multivariate code, is outside the obstruction.

\section{What the quadratic-capacity construction actually uses}

Koiran and Sontag, \emph{Neural Networks with Quadratic VC Dimension},
JCSS 54 (1997), 190--198,
\href{https://sontaglab.org/FTPDIR/quadratic-vc.pdf}{author-hosted paper},
Proposition 1 and Theorems 2, 4, and 6, give two mechanisms.
The first selects a stored real, repeatedly extracts a bit, updates
the remainder, and finally selects an output bit. Its sigmoid
simulation uses large weights for threshold gates and small weights
for linear gates, preserving the sequential composition. The second
uses iterates of $\phi(w)=4w(1-w)$, whose symbolic itineraries encode
arbitrary words. Input and output selection use Lagrange interpolation,
written as a product times a sum of reciprocals. Those reciprocal
denominators have interior signed poles. The more general smooth
activation result simulates multiplication by a second-order Taylor
limit. None of these constructions is a single hidden layer of
exterior affine reciprocals.

The following statements isolate a precise cost of flattening the
scalar decoder. They apply even when the coefficients of each
individual query are otherwise unrestricted.

\section{An exact oscillation obstruction}

\begin{lemma}[Degree after clearing positive affine denominators]
Let
\[
 R(w)=b+\sum_{h=1}^H\frac{a_h+c_hw}{d_h+e_hw},
 \qquad d_h+e_hw>0\quad(0\le w\le1).
\]
If $R$ has alternating strict signs at ordered points
$w_0<\cdots<w_m$ in $[0,1]$, then $H\ge m$.
\end{lemma}
\begin{proof}
Multiplication by the positive product of the denominators preserves
every sign and yields the polynomial
\[
 P(w)=b\prod_h(d_h+e_hw)
       +\sum_h(a_h+c_hw)\prod_{g\ne h}(d_g+e_gw),
\]
of degree at most $H$. It is not the zero polynomial because its
values at the specified points are nonzero. The intermediate value
theorem gives a different zero in each of the $m$ disjoint intervals
$(w_{j-1},w_j)$. Thus $\deg P\ge m$.
\end{proof}

\begin{theorem}[Flattening a logistic iterate costs exponentially]
For every integer $k\ge0$, a sum of $H$ affine ratios with
denominators positive on $[0,1]$ that agrees in strict sign with
$\phi^{\circ k}(w)-1/2$ at all of its successive extrema must have
$H\ge2^k$.
\end{theorem}
\begin{proof}
The elementary double-angle identity gives
\[
 \phi(\sin^2t)=\sin^2(2t),\qquad
 \phi^{\circ k}(\sin^2t)=\sin^2(2^kt).
\]
Put
\[
 w_j=\sin^2\left(\frac{j\pi}{2^{k+1}}\right),
             \qquad 0\le j\le2^k.
\]
These points increase strictly from zero to one, and
$\phi^{\circ k}(w_j)=\sin^2(j\pi/2)$ alternates between zero
and one. Apply the preceding lemma to the $2^k$ resulting sign
changes. The proof uses only this finite set of test points.
\end{proof}

This theorem permits arbitrary cancellation among all heads. Its
hypothesis also follows if each affine denominator is positive only
at the two endpoints, since an affine function positive at both
endpoints is positive throughout the interval. It therefore cannot
be evaded by putting a pole between the test points while retaining
endpoint admissibility.

For a decoder with $q$ successive logistic itinerary bits, the last
bit uses $k=q-1$, so the same direct flat realization needs at least
$2^{q-1}$ heads. This is a restriction on that decoder, not a claim
that every possible $q$-bit code has its particular itinerary order.

\section{Arbitrary scalar codes still require many scalar crossings}

\begin{theorem}[A scalar-reader pattern bound]
Suppose $q$ fixed query scores $R_1(w),\ldots,R_q(w)$ each have
the form in the lemma with at most $H$ heads, and every denominator
is positive on one common interval. If varying the scalar memory
$w$ within that interval realizes all $2^q$ strict sign words, then
\[
                         2^q\le qH+1.
\]
\end{theorem}
\begin{proof}
Clear the positive denominators separately in each query, obtaining
polynomials $P_i$ of degree at most $H$. An identically zero $P_i$
would prevent strict shattering, so no $P_i$ is zero. Their combined
real zero set in the interval has at most $qH$ points. On each of
its at most $qH+1$ complementary intervals, all $q$ signs are
constant. At a zero at least one sign is not strict. Thus at most
$qH+1$ strict words can occur.
\end{proof}

The theorem permits query-specific affine coefficients and poles.
It does not cover a memory map $w\mapsto\theta(w)$ whose coordinates
are themselves unrestricted nonlinear functions of $w$. Such a
map must be supplied and audited as a new construction; it is not
part of a gate-by-gate transfer of the scalar decoder.

\section{The local operation that does transfer}

There is no obstruction to approximating one multiplication of
affine input functions. For example, let $L$ be any fixed affine
function on a finite Boolean cube, let
$D_t=1+\sqrt t\sum_i x_i$, and take $t>0$ sufficiently small.
Then both $D_t\pm tL$ are positive and have strictly positive
slopes, and
\[
 \frac1{2t}\left(\frac{L}{D_t-tL}-\frac{L}{D_t+tL}\right)
       =\frac{L^2}{D_t^2-t^2L^2}\longrightarrow L^2
\]
uniformly on the cube. Thus two admissible heads give an affine
square in the limit. Polarization gives products of affine
functions at constant cost. A later gate in the logistic decoder,
however, receives a nonlinear polynomial, not an affine input.
The exponential crossing theorem proves that recursively reusing
this local identity cannot be flattened at constant cost per
iteration.

The interpolation reciprocals in Theorem 6 of the cited paper
also fail a direct transfer: $j-\epsilon-i$ takes both signs as
$j$ ranges over the interpolation nodes whenever $i$ is interior.
No single sign change of that denominator makes it positive at
all nodes. The accompanying product is an essential part of the
value interpolation and is not an admissible affine numerator.

\section{Smooth parameter dimension does not create a logarithm}

Lee, Bartlett, and Williamson's 1997
\href{https://doi.org/10.1162/neco.1997.9.4.765}{published correction}
withdraws the generic boundary-uniqueness theorem used in their 1995
article. The correction retains a lower bound proportional to the
number of parameters for the tanh case, including small parameter
ranges. The
\href{https://www.researchgate.net/publication/242919186_Correction_to_Lower_Bounds_on_VC-Dimension_of_Smoothly_Parameterized_Function_Classes}{publisher preview}
states the correction explicitly. It supplies no extra logarithmic
factor and no flat rational bit decoder.

A direct counterexample to the unqualified uniqueness principle is
$F(a,x)=\|x-a\|^2$ with $a,x\in\mathbb R^d$, $d\ge2$. Its zero
sets are distinct singletons for distinct $a$, but the classifiers
$\mathbf1[F(a,x)>0]$ have VC dimension one: a singleton can receive
either label, while two distinct points cannot both receive zero.
The parameter manifold has dimension $d$. Under the stricter
convention that all classified values must be nonzero, the same
family has only the all-positive strict pattern.

For comparison, the valid elementary local criterion is this:
if the evaluation map
\[
 a\longmapsto(F(a,x_1),\ldots,F(a,x_m))
\]
vanishes at some parameter and has differential of rank $m$ there,
then its image contains a neighborhood of zero, by restricting to
an $m$-dimensional coordinate slice and applying the inverse
function theorem. Hence it realizes all $2^m$ strict orthants
locally. This criterion requires $m$ no larger than the parameter
dimension. It cannot by itself prove the desired superlinear
information per real parameter.

\paragraph{Conclusion of the transfer audit.}
The deep constructions are unconditional memory constructions,
and their one-step Taylor simulation is compatible with the
original affine-ratio model. Their repeated decoder and their
interior interpolation poles do not transfer to one layer at
comparable head cost. The missing step is a different joint memory
decoder, not an omitted routine depth reduction. No global upper
bound improvement follows from these sources.
\end{document}
